% ===== PRÉAMBULE CANONIQUE — à coller VERBATIM en tête de CHAQUE .tex =====
\documentclass[11pt,a4paper]{article}
\usepackage[utf8]{inputenc}\usepackage[T1]{fontenc}\usepackage{lmodern}
\usepackage[french]{babel}
\usepackage{amsmath,amssymb,mathtools}\usepackage{icomma}
\usepackage[version=4]{mhchem}          % \ce{...} : équations & formules
\usepackage{siunitx}\sisetup{locale=FR} % \qty{}{}, \si{}, \ang{}
\usepackage{chemfig}                    % \chemfig{...} : structures développées
\usepackage{xcolor}\usepackage{colortbl}
\usepackage{tikz}\usepackage{pgfplots}\pgfplotsset{compat=1.18}
\usepackage[most]{tcolorbox}\usepackage{enumitem}\usepackage{array,booktabs}
\usepackage[a4paper,margin=2.2cm]{geometry}
\usetikzlibrary{arrows.meta,calc,angles,quotes,positioning,patterns,backgrounds,shapes.geometric}
\definecolor{primary}{RGB}{222,49,99}\definecolor{primarydark}{RGB}{150,22,55}
\definecolor{defcol}{RGB}{0,114,178}\definecolor{methcol}{RGB}{52,92,148}
\definecolor{excol}{RGB}{0,135,90}\definecolor{attncol}{RGB}{200,40,55}
\tcbset{boxbase/.style={enhanced,breakable,boxrule=0.9pt,arc=2.5pt,
  left=3mm,right=3mm,top=2.3mm,bottom=2.3mm,fonttitle=\bfseries,coltitle=white}}
% --- boîtes TUTORIEL (noms EXACTS, ne pas renommer) ---
\newtcolorbox{cle}[1][]{boxbase,colback=defcol!6,colframe=defcol,
  title={\textsc{À retenir}\ifx\relax#1\relax\relax\else\textmd{ : \itshape #1}\fi}}
\newtcolorbox{methode}[1][]{boxbase,colback=methcol!7,colframe=methcol,sharp corners,
  title={\textsc{Méthode}\ifx\relax#1\relax\relax\else\textmd{ --- \itshape #1}\fi}}
\newtcolorbox{exemplebox}[1][]{boxbase,colback=excol!5,colframe=excol,
  title={\textsc{Exemple}\ifx\relax#1\relax\relax\else\textmd{ : \itshape #1}\fi}}
\newtcolorbox{piege}[1][]{boxbase,colback=attncol!6,colframe=attncol,
  title={\textsc{Piège}\ifx\relax#1\relax\relax\else\textmd{ : \itshape #1}\fi}}
\newtcolorbox{remarque}[1][]{boxbase,colback=black!4,colframe=black!45,
  title={\textsc{Remarque}\ifx\relax#1\relax\relax\else\textmd{ : \itshape #1}\fi}}
% --- boîtes EXERCICE / CORRIGÉ (noms EXACTS, auto-comptées) ---
\newtcolorbox[auto counter]{exo}[1][]{boxbase,colback=primary!4,colframe=primary,
  title={\textsc{Exercice}~\thetcbcounter\ifx\relax#1\relax\relax\else\textmd{ --- \itshape #1}\fi}}
\newtcolorbox[auto counter]{corrige}[1][]{boxbase,colback=excol!4,colframe=excol,
  title={\textsc{Corrigé}~\thetcbcounter\ifx\relax#1\relax\relax\else\textmd{ --- \itshape #1}\fi}}
\newcommand{\R}{\mathbb{R}}\newcommand{\N}{\mathbb{N}}\newcommand{\Z}{\mathbb{Z}}\newcommand{\ds}{\displaystyle}
% =========================================================================
\begin{document}
\begin{exo}[chaleur sensible $Q=mc\Delta T$]
On donne $c_{\text{eau}}=\qty{4.185}{\kilo\joule\per\kelvin\per\kilogram}$. Calcule la chaleur $Q$ à fournir.
\begin{enumerate}[label=\alph*)]
  \item chauffer $\qty{2.0}{\kilogram}$ d'eau de $\qty{20}{\celsius}$ à $\qty{30}{\celsius}$ ;
  \item chauffer $\qty{0.50}{\kilogram}$ d'eau de $\qty{15}{\celsius}$ à $\qty{100}{\celsius}$.
\end{enumerate}
\end{exo}

\begin{exo}[chaleur latente $Q=mL$]
On donne $L_{\text{fusion}}=\qty{334}{\kilo\joule\per\kilogram}$ et $L_{\text{vaporisation}}=\qty{2257}{\kilo\joule\per\kilogram}$. Calcule $Q$.
\begin{enumerate}[label=\alph*)]
  \item faire fondre $\qty{0.20}{\kilogram}$ de glace à $\qty{0}{\celsius}$ ;
  \item vaporiser $\qty{0.10}{\kilogram}$ d'eau à $\qty{100}{\celsius}$.
\end{enumerate}
\end{exo}

\begin{exo}[chaleur reçue ou cédée ?]
Indique le signe de $Q$ (reçue $>0$ ou cédée $<0$), ou la relation demandée.
\begin{enumerate}[label=\alph*)]
  \item un corps passe de $\qty{20}{\celsius}$ à $\qty{60}{\celsius}$ ;
  \item un corps passe de $\qty{80}{\celsius}$ à $\qty{25}{\celsius}$ ;
  \item deux corps seuls dans un calorimètre adiabatique : quelle relation lie leurs chaleurs échangées ?
\end{enumerate}
\end{exo}

\begin{exo}[conversions indispensables]
Convertis dans l'unité demandée.
\begin{enumerate}[label=\alph*)]
  \item $\qty{4.185}{\kilo\joule}$ en $\si{\joule}$ ;
  \item $\qty{250}{\gram}$ en $\si{\kilogram}$ ;
  \item une variation de température de $\qty{25}{\celsius}$ en $\si{\kelvin}$.
\end{enumerate}
\end{exo}

\begin{exo}[comparer deux dépenses d'énergie]
On dispose de $\qty{0.10}{\kilogram}$ d'eau. Données : $c_{\text{eau}}=\qty{4.185}{\kilo\joule\per\kelvin\per\kilogram}$, $L_{\text{fusion}}=\qty{334}{\kilo\joule\per\kilogram}$.
\begin{enumerate}[label=\alph*)]
  \item chaleur pour \emph{faire fondre} $\qty{0.10}{\kilogram}$ de glace à $\qty{0}{\celsius}$ ;
  \item chaleur pour \emph{chauffer} $\qty{0.10}{\kilogram}$ d'eau liquide de $\qty{0}{\celsius}$ à $\qty{50}{\celsius}$ ;
  \item laquelle des deux opérations coûte le plus d'énergie ?
\end{enumerate}
\end{exo}

\end{document}
